\documentclass[12pt,a4paper]{article}
\usepackage{fontspec}
\usepackage{polyglossia}
\setdefaultlanguage{ukrainian}
\setmainfont{NewComputerModern10}
\setsansfont{Noto Sans}
\setmonofont{Noto Sans Mono}[
  Scale=.82,
  ItalicFont={Noto Sans Mono},
  BoldItalicFont={Noto Sans Mono Bold}
]
\usepackage{amsmath,amssymb,booktabs,caption,enumitem,float,geometry,graphicx,hyperref,microtype,multirow,titlesec,xcolor}
\geometry{left=2cm,right=2cm,top=2cm,bottom=2cm}
\setlength{\parindent}{1.25cm}
\linespread{1.08}
\titleformat{\section}{\normalfont\fontsize{14}{16}\selectfont\bfseries}{\thesection}{.6em}{}
\titleformat{\subsection}{\normalfont\fontsize{12}{14}\selectfont\bfseries}{\thesubsection}{.6em}{}
\titlespacing*{\section}{0pt}{1.3em}{.6em}
\captionsetup{font=small,labelfont=bf,labelsep=endash}
\usepackage[newfloat]{minted}
\SetupFloatingEnvironment{listing}{name=Лістинг}
\setminted{fontsize=\footnotesize,linenos,breaklines,autogobble,frame=single,framesep=2mm,bgcolor=gray!5,tabsize=4}
\hypersetup{unicode,colorlinks=true,linkcolor=black,pdftitle={ІО-41, Давидчук А. М., Модульна контрольна робота},pdfauthor={Давидчук А. М.}}
\newcommand{\resultDir}{../results}
\newcommand{\srcDir}{../src}
\newcommand{\studentGroup}{ІО-41}
\newcommand{\studentName}{Давидчук А. М.}
\newcommand{\reportYear}{2026}

\begin{document}
\pagenumbering{gobble}
\begin{titlepage}
\begin{center}\includegraphics[width=.7\textwidth]{../../../../../templates/letterhead.png}\end{center}
\begin{center}
\large Національний технічний університет України\\
«Київський політехнічний інститут імені Ігоря Сікорського»\\[1em]
Факультет інформатики та обчислювальної техніки\\
Кафедра обчислювальної техніки
\end{center}
\vfill
\begin{center}
\textbf{\LARGE Методи та технології штучного інтелекту}\\[2em]
\textbf{\Large Модульна контрольна робота}\\[.7em]
\parbox{.92\textwidth}{\centering «Навчання нейронної мережі прямого поширення генетичним алгоритмом»}
\end{center}
\vfill
\begin{flushright}
\begin{minipage}{.48\textwidth}
Виконав: студент 3 курсу ФІОТ,\\
групи \studentGroup\\
\textit{\studentName}\\[1em]
Перевірив: \textit{Польшакова О. М.}
\end{minipage}
\end{flushright}
\vfill
\begin{center}Київ --- \reportYear\end{center}
\end{titlepage}
\pagenumbering{arabic}

\section{Завдання та індивідуальний варіант}
Необхідно реалізувати нейронну мережу прямого поширення та навчити її генетичним алгоритмом моделювати функцію двох змінних. Для студента Давидчука Артема Миколайовича у таблиці варіантів задано структуру
\begin{equation}
2-4-8-10-6-6-1
\end{equation}
і функцію
\begin{equation}
z=x\sin|y|.
\end{equation}

Крім основного завдання, реалізовано обидва розширення на 10 балів: модифіковану ініціалізацію Нгуєна--Відроу з урахуванням значущості ознак та власне зворотне поширення помилки з оптимізатором Adam. Також виконано розширення на 20 балів: GUI-програму для автоматичного структурного синтезу мережі генетичним алгоритмом.

В умові не визначено конкретний інтервал для варіанта 52. Для відтворюваного експерименту прийнято $x,y\in[0;1]$. Межі винесені в параметри програми й можуть бути змінені без перероблення алгоритмів.

\section{Математична модель нейронної мережі}
Мережа має два входи, п'ять прихованих шарів із 4, 8, 10, 6 і 6 нейронами та один вихід. Для прихованих шарів використано гіперболічний тангенс
\begin{equation}
a^{(l)}=\tanh\left(a^{(l-1)}W^{(l)}+b^{(l)}\right),
\end{equation}
а вихідний шар є лінійним:
\begin{equation}
\hat z=a^{(5)}W^{(6)}+b^{(6)}.
\end{equation}
Лінійний вихід не обмежує діапазон прогнозу й відповідає задачі регресії.

\begin{figure}[H]
\centering
\includegraphics[width=.98\textwidth]{\resultDir/architecture.png}
\caption{Задана архітектура нейронної мережі}
\end{figure}

Кількість параметрів шару з $n_{l-1}$ входами та $n_l$ виходами дорівнює $n_{l-1}n_l+n_l$. Для заданої структури загальна кількість ваг і зміщень:
\begin{equation}
(2\cdot4+4)+(4\cdot8+8)+(8\cdot10+10)+(10\cdot6+6)+(6\cdot6+6)+(6\cdot1+1)=257.
\end{equation}
Тому хромосома генетичного алгоритму містить 257 дійсних генів.

\section{Генетичне навчання мережі}
Кожна особина популяції є повним набором параметрів мережі. Пристосованість оцінюється середньоквадратичною похибкою
\begin{equation}
\mathrm{MSE}=\frac{1}{N}\sum_{i=1}^{N}(z_i-\hat z_i)^2,
\end{equation}
яку потрібно мінімізувати.

Алгоритм складається з таких кроків:
\begin{enumerate}[leftmargin=*]
\item створення початкової популяції хромосом;
\item пряме поширення для кожної хромосоми й розрахунок MSE;
\item турнірний відбір батьків;
\item арифметичний blend-кросовер;
\item гаусова мутація окремих генів зі спадним масштабом;
\item перенесення чотирьох найкращих особин без змін;
\item повторення до 280 поколінь або спрацювання критерію відсутності покращення.
\end{enumerate}

\begin{listing}[H]
\inputminted[firstline=184,lastline=230]{python}{\srcDir/mkr_core.py}
\caption{Оцінювання, відбір, кросовер і мутація}
\end{listing}

\section{Модифікована ініціалізація Нгуєна--Відроу}
Звичайна випадкова ініціалізація не використовує інформацію про залежність виходу від входів. У модифікованому варіанті спочатку оцінюється значущість кожної ознаки через абсолютну кореляцію зі значенням $z$:
\begin{equation}
s_j=\frac{|\operatorname{corr}(x_j,z)|+\varepsilon}
{\sum_k(|\operatorname{corr}(x_k,z)|+\varepsilon)}.
\end{equation}
Для сформованої вибірки отримано $s_x=0{,}5218$ і $s_y=0{,}4782$. Стовпці ваг першого шару масштабуються відповідно до $\sqrt{2s_j}$, після чого для кожного шару застосовується нормування Нгуєна--Відроу з коефіцієнтом
\begin{equation}
\beta=0{,}7n_{\mathrm{out}}^{1/n_{\mathrm{in}}}.
\end{equation}

\begin{listing}[H]
\inputminted[firstline=114,lastline=144]{python}{\srcDir/mkr_core.py}
\caption{Оцінювання значущості ознак та модифікована ініціалізація}
\end{listing}

\section{Самостійна реалізація backpropagation та Adam}
Для порівняння реалізовано зворотне поширення помилки без готової функції машинного навчання. Для вихідного шару градієнт дорівнює
\begin{equation}
\delta^{(L)}=\frac{2}{N}(\hat z-z),
\end{equation}
а для прихованих шарів
\begin{equation}
\delta^{(l)}=\left(\delta^{(l+1)}(W^{(l+1)})^T\right)\odot\left(1-(a^{(l)})^2\right).
\end{equation}
Градієнти ваг і зміщень обчислюються як $\nabla W^{(l)}=(a^{(l-1)})^T\delta^{(l)}$ та $\nabla b^{(l)}=\sum_i\delta_i^{(l)}$. Adam накопичує перший і другий моменти градієнта та виконує корекцію зміщення моментів.

\begin{listing}[H]
\inputminted[firstline=267,lastline=296]{python}{\srcDir/mkr_core.py}
\caption{Власна реалізація backpropagation та Adam}
\end{listing}

\section{Програмна реалізація та параметри експерименту}
Вибірку сформовано на сітці $31\times31$, тобто отримано 961 точку. Випадковим, але відтворюваним розбиттям 769 точок використано для навчання, а 192 --- для незалежного тестування. У прихованих шарах використано $\tanh$, на виході --- лінійну активацію.

\begin{table}[H]
\centering
\begin{tabular}{lll}
\toprule
Метод & Параметр & Значення\\
\midrule
\multirow{7}{*}{Генетичний алгоритм}
& Розмір популяції & 80\\
& Максимум поколінь & 280\\
& Елітних особин & 4\\
& Розмір турніру & 4\\
& Імовірність кросоверу & 0,90\\
& Імовірність мутації гена & 0,035\\
& Масштаб мутації & 0,12 $\rightarrow$ 0,008\\
\midrule
\multirow{3}{*}{Adam}
& Кількість епох & 1200\\
& Розмір пакета & 64\\
& Швидкість навчання & 0,003\\
\bottomrule
\end{tabular}
\caption{Параметри навчання}
\end{table}

Клас мережі не залежить від конкретної кількості шарів: архітектура передається списком, а ваги та зміщення автоматично пакуються в одну хромосому.

\begin{listing}[H]
\inputminted[firstline=48,lastline=73]{python}{\srcDir/mkr_core.py}
\caption{Опис архітектури та розпакування хромосоми}
\end{listing}

\begin{listing}[H]
\inputminted[firstline=75,lastline=92]{python}{\srcDir/mkr_core.py}
\caption{Пакування параметрів і пряме поширення сигналу}
\end{listing}

\section{Результати навчання заданої мережі}
У таблиці наведено метрики на тестовій вибірці й час одного запуску на використаному комп'ютері. Час залежить від обладнання, але дає змогу порівняти алгоритми в однакових умовах.

\begin{table}[H]
\centering
\small
\begin{tabular}{lrrrrr}
\toprule
Метод & MSE & RMSE & MAE & $\max|e|$ & Час, с\\
\midrule
ГА, випадкова ініціалізація & 0,00115422 & 0,033974 & 0,028288 & 0,088683 & 10,14\\
ГА, Нгуєн--Відроу & 0,00012166 & 0,011030 & 0,008009 & 0,041776 & 10,03\\
Backpropagation, Adam & 0,00007381 & 0,008592 & 0,007133 & 0,018942 & 2,77\\
\bottomrule
\end{tabular}
\caption{Порівняння методів навчання на тестовій вибірці}
\end{table}

Модифікована ініціалізація зменшила тестову MSE генетичного алгоритму у 9,49 раза порівняно з випадковою. Отже, використання апріорної інформації та геометрично узгодженого масштабу початкових ваг суттєво покращило пошук у 257-вимірному просторі.

Adam досяг MSE $7{,}38\cdot10^{-5}$, що у 1,65 раза менше за результат найкращого ГА. Крім того, його час був приблизно у 3,6 раза меншим. Однак backpropagation потребує диференційовної моделі та явного виведення градієнтів, тоді як ГА працює лише зі значенням функції якості та простіше адаптується до дискретних або недиференційовних параметрів.

\begin{figure}[H]
\centering
\includegraphics[width=.98\textwidth]{\resultDir/convergence.png}
\caption{Збіжність генетичного алгоритму та Adam}
\end{figure}

Криві ГА зменшуються монотонно завдяки елітизму. Ініціалізація Нгуєна--Відроу після приблизно 30 поколінь випереджає випадкову й завершує пошук на порядок нижче. Крива Adam має коливання через мініпакетне навчання, але загальна тенденція також спадна.

\begin{figure}[H]
\centering
\begin{minipage}[t]{.49\textwidth}\centering
\includegraphics[width=\textwidth]{\resultDir/target_surface.png}
\captionof{figure}{Точна поверхня $z=x\sin|y|$}
\end{minipage}\hfill
\begin{minipage}[t]{.49\textwidth}\centering
\includegraphics[width=\textwidth]{\resultDir/prediction_surface.png}
\captionof{figure}{Прогноз після навчання ГА}
\end{minipage}
\end{figure}

Прогноз відтворює нульові значення при $x=0$ або $y=0$ та плавне зростання до $\sin(1)\approx0{,}8415$ у точці $(1;1)$.

\begin{figure}[H]
\centering
\begin{minipage}[t]{.49\textwidth}\centering
\includegraphics[width=\textwidth]{\resultDir/error_surface.png}
\captionof{figure}{Поверхня абсолютної похибки}
\end{minipage}\hfill
\begin{minipage}[t]{.49\textwidth}\centering
\includegraphics[width=\textwidth]{\resultDir/slice_comparison.png}
\captionof{figure}{Переріз при $y=0{,}5$}
\end{minipage}
\end{figure}

Найбільші відхилення зосереджені біля меж області, але максимальна тестова похибка не перевищує 0,0418. На контрольному перерізі прогноз практично повторює точну функцію.

\begin{figure}[H]
\centering
\includegraphics[width=.72\textwidth]{\resultDir/method_comparison.png}
\caption{Порівняння тестової MSE}
\end{figure}

\section{GUI та генетичний структурний синтез}
GUI приймає математичний вираз, інтервали $x$ і $y$, максимальну сумарну кількість прихованих нейронів, допустиму MSE та параметри зовнішнього ГА. Вираз розбирається через модуль \texttt{ast}; дозволені лише змінні $x$, $y$, арифметичні операції та визначений набір математичних функцій.

Структура кодується вектором із п'яти цілих генів. Нуль означає відсутній шар, а додатне число --- кількість нейронів. Зовнішній ГА змінює структуру, а внутрішній ГА навчає ваги кандидата. Функція пристосованості містить валідаційну MSE і малий штраф за кількість нейронів:
\begin{equation}
F(A)=\mathrm{MSE}_{\mathrm{val}}(A)+2\cdot10^{-5}\sum_{l}n_l.
\end{equation}

\begin{listing}[H]
\inputminted[firstline=69,lastline=94]{python}{\srcDir/structure_search.py}
\caption{Оцінювання архітектури внутрішнім генетичним алгоритмом}
\end{listing}

\begin{listing}[H]
\inputminted[firstline=96,lastline=129]{python}{\srcDir/structure_search.py}
\caption{Еволюція структур нейронної мережі}
\end{listing}

Пошук виконується в окремому потоці, щоб вікно не блокувалося. У додатковому вікні відображаються параметри ГА та найкраща проміжна структура кожного покоління.

\begin{figure}[H]
\centering
\includegraphics[width=.80\textwidth]{\resultDir/gui_preview.png}
\caption{Схематичний вигляд GUI та його результатів}
\end{figure}

У контрольному headless-запуску обмежено сумарну кількість прихованих нейронів значенням 24. За шість поколінь оцінено 12 унікальних архітектур. Найкращою стала мережа $2-16-1$ з валідаційною MSE 0,0019517.

\begin{figure}[H]
\centering
\includegraphics[width=.60\textwidth]{\resultDir/structure_search.png}
\caption{Зміна найкращої похибки під час структурного синтезу}
\end{figure}

\section{Висновки}
Реалізовано задану нейронну мережу $2-4-8-10-6-6-1$ для моделювання функції $z=x\sin|y|$. Усі 257 ваг і зміщень закодовано однією дійсною хромосомою та навчено генетичним алгоритмом. За модифікованої ініціалізації Нгуєна--Відроу тестова MSE становила 0,00012166, що у 9,49 раза краще за випадкову ініціалізацію.

Самостійно реалізований backpropagation з Adam отримав тестову MSE 0,00007381 та потребував 2,77 с проти приблизно 10 с для кожного запуску ГА. Adam виявився ефективнішим для цієї гладкої диференційовної функції, а ГА продемонстрував незалежність від градієнтів і придатність до одночасної оптимізації структури.

Реалізовано GUI для структурного синтезу. Програма приймає функцію та обмеження, показує параметри й проміжні результати ГА та повертає структуру мережі з її похибкою. Таким чином виконано основне завдання та всі три запропоновані розширення.

\end{document}
